% !PRINTED 236
\weberpage{236}

%<<CONT
identical substitution, then we must have $\varphi^2 = 1$, and this is also sufficient. Here we can, by transformation, bring $\varphi$ to the form $\begin{pmatrix} \alpha, & 0 \\ 0, & -\alpha^{-1} \end{pmatrix}$, and obtain as condition $\alpha^2 = \pm 1$, and accordingly two forms of the substitution $\varphi$:
\[
\varphi' = \begin{pmatrix} i, & 0 \\ 0, & i \end{pmatrix}, \quad
\varphi'' = \begin{pmatrix} 1, & 0 \\ 0, & -1 \end{pmatrix},
\]
of which the first is called the \gesperrt{inversion}, the second the \gesperrt{reflection}. If one carries them over, according to \S.\,50, (1), to a rectangular coordinate transformation, then $\varphi'$ signifies the simultaneous change of sign of all three coordinates, $\varphi''$ the replacement of $x$ by $-x$. From this there result the two groups of the second kind:
\weq{2}{C_1' = \begin{pmatrix} i^h, & 0 \\ 0, & i^h \end{pmatrix}, \quad C_1'' = \begin{pmatrix} 1, & 0 \\ 0, & (-1)^h \end{pmatrix}, \quad h = 0, 1.}

II. Let $P$ be a cyclic group $C_n$, $n > 1$.
If we set
\weq{3}{c = \begin{pmatrix} \varepsilon, & 0 \\ 0, & \varepsilon^{-1} \end{pmatrix}, \quad \varepsilon = e^{\frac{\pi i}{n}},}
then by \S.\,55:
\weq{4}{C_n = 1, c, c^2, \ldots, c^{n-1},}
in which
\[
c^s = \begin{pmatrix} \varepsilon^s, & 0 \\ 0, & \varepsilon^{-s} \end{pmatrix};
\]
if now
\weq{5}{\varphi = \begin{pmatrix} \alpha, & \beta \\ \gamma, & \delta \end{pmatrix}, \quad \alpha\delta - \beta\gamma = -1}
holds, then we first form
\weq{6}{\varphi^{-1} c \varphi = \begin{pmatrix} \alpha\delta\varepsilon - \beta\gamma\varepsilon^{-1}, & \beta\delta(\varepsilon - \varepsilon^{-1}) \\ -\alpha\gamma(\varepsilon - \varepsilon^{-1}), & \alpha\delta\varepsilon^{-1} - \beta\gamma\varepsilon \end{pmatrix},}
\weq{7}{\varphi^2 = \begin{pmatrix} \alpha^2 + \beta\gamma, & (\alpha + \delta)\beta \\ (\alpha + \delta)\gamma, & \beta\gamma + \delta^2 \end{pmatrix}.}
These two substitutions must occur in the sequence (4), and hence we must have $\beta\delta = 0$, $\alpha\gamma = 0$; thus either $\beta$ and $\gamma$ are both zero, or $\alpha$ is zero together with $\delta = 0$. If $\beta = \gamma = 0$, it follows from (7) that $\alpha$ is of the form $\varepsilon^{1/2\,r}$, and hence
\[
\varphi = \begin{pmatrix} e^{\frac{\pi i}{2n}\,r}, & 0 \\ 0, & -e^{-\frac{\pi i}{2n}\,r} \end{pmatrix}
\]
%CONT>>
